\section{The actual preordinary classes and their quadratic ray norms}
\label{sec:preordinary-ray-classes}

This section uses AC-K33-ray-norm.  It concerns the actual
Rankin--Iwasawa classes, rather than an abstract collection of
coefficient maps.  We take \(M_0=1\) in that theorem and choose
\(N_0\geq5\) divisible by \(3\) and by the tame level of the
fixed eigenform.  The auxiliary ray ideals are prime to this
fixed level, to \(p=269\), and to one fixed smoothing integer
\(c>1\) coprime to \(6N_0p\).  We work at cyclotomic modulus \(m=1\); thus there is
no cyclotomic \(\sigma_q\) factor to suppress.

Let \(f\) be the actual cuspidal eigenform of weight \(2e+1\)
and quadratic nebentypus \(\eta_K\) occurring in AC-K33-image,
put \(k=2e-1\) and \(j=e-1\), and keep its declared coefficient
normalization
\[
 M(f)^*=\rho_f,\qquad
 \det(1-\operatorname{Frob}_q^{\rm arith}X\mid\rho_f)
       =1-a_q(f)X+q^{2e}\eta_K(q)X^2
 \quad(q\nmid N_fp).
 \tag{RCN.1}
\]
The target over \(K\) is
\[
 T=\rho_f|_{G_K}\otimes\chi^{-b}(-j).
 \tag{RCN.2}
\]
By the explicitly chosen determinant/intertwiner identification
in AC-K33-image, this is the actual coefficient
\((V|_{G_K}\chi^b)^\vee(1)\).  Let
\[
 P_{\mathfrak l}(X)
   =\det\bigl(1-
       \operatorname{Frob}_{\mathfrak l}^{\rm geom}X
                    \mid T^\vee(1)\bigr).
 \tag{RCN.3}
\]
All primes considered below are unramified in this fixed
coefficient.

The primary geometric identities used in the proof are exactly
KLZ, arXiv:1503.02888v6, Theorem 5.3.1, Theorem 5.5.1,
Corollary 5.5.2 and (5.6.1a)--(5.6.1b), printed pp.~25--28.
Definition 5.1.5 is the actual preordinary class; the scalar
\([q]_*\) acts only on the measure factor by Notation 5.2.1.
The calculations below prove the needed asymmetric consequences
of these identities.  They do not infer them by substituting a
negative weight into the classical-eigenform statement 5.7.6.
The proof of Theorem 5.5.1 in the source refers to its geometric
antecedent LLZ15, Lemma A.2.1; that published geometric identity
is retained here; its complete printed proof, including both
correspondence diagrams and the diamond conversion, has been read.
No new reproof of the elliptic polylogarithm or its norm relation
is claimed.

\begin{lemma}[AC-K33-ray-Hecke: the combined torsor scalar]
\label{AC-K33-ray-Hecke}
At a rational prime \(q\) unramified in the second coefficient,
let \(a,b\in\overline{A_{\mathfrak n}[1/p]}\) be its geometric
Frobenius roots, and put \(Q=q^{j+1}\).
On the actual quotient \(Q_{\mathfrak n}\), the combined
operator of KLZ Notation 5.2.1 has scalar
\[
 K_g=[q]_*\langle q^{-1}\rangle=\frac{ab}{Q}.
 \tag{RCN.4}
\]
This is an identity over \(A_{\mathfrak n}\); adjoining the
roots is unnecessary.  The analogous first-factor scalar is
\[
 K_f=\frac{d_f}{Q},\qquad
 t_f=a_q(f),\qquad d_f=q^{2e}\eta_K(q).
 \tag{RCN.5}
\]
If the second local profile is \(v_b\) of (RNN.6), then
\(U_q'=b\) on its quotient; for the compact profile it is zero.
\end{lemma}
\begin{proof}
We derive (RCN.4) from the actual correspondence identities,
thereby keeping the measure and diamond conventions together.
Let \(i=\operatorname{pr}_1^*\) on the compact dual tower.
The transpose of (5.6.1a) is
\[
 U_q i=iT_q-\operatorname{pr}_2^*K_g^t .
 \tag{RCN.6}
\]
The operators here are the source operators, dual to the
primed ones.  In the standard local action \(U_q\) is
the sum over
\(\left(\begin{smallmatrix}q&r\\0&1\end{smallmatrix}\right)\),
\(0\leq r<q\).  On a valuation-dependent Kirillov function
supported on \(\mathbb Z_q\), it is \(qv(qx)\) for
\(x\in\mathbb Z_q\); for \(v_q(x)=-1\), the sum of the
nontrivial additive \(q\)-th roots is zero, and for smaller
valuation the support is zero.
Thus (RNN.5) gives the exact identity
\[
 U_q v_{a,b}=(a+b)v_{a,b}-\frac{ab}{q}D_qv_{a,b}.
 \tag{RCN.7}
\]
The spherical Hecke operator is \(T_q=a+b\), in the source
normalization already checked in AC-K33-ray-refinement.
By (RNN.11), \(\operatorname{pr}_2^*=q^jD(q)\), and the
non-\(q\) unit components fix our tensor.
Substitution into (RCN.6), followed by cancellation of the
invertible group action \(D(q)\), proves (RCN.4) on the actual
compact cohomology vector.  Taking the bounded dual proves
the claimed quotient identity.
In particular we have not assigned a separate, guessed
character to \([q]_*\).

For the first factor, the literal moment
\(\mathrm{mom}^{k-j}\cdot\mathrm{id}\) gives
\([q]_*\langle q^{-1}\rangle=q^{k-j}\eta_K(q)\),
as in the matrix calculation in the complete proof of
KLZ Theorem 5.7.6.  Since \(k+1=2e\), this is (RCN.5).
This calculation uses the allowed positive moment \(k-j=e\).
Finally the same additive-root sum applied to (RNN.6)
gives \(qv_b(qx)=b v_b(x)\); applied to
\(\delta_q\) it gives zero.  These identities commute with
all coefficient reductions and are polynomial in the
unramified Frobenius values.
\end{proof}

\begin{theorem}[AC-K33-ray-classes: quadratic norms of the actual images]
\label{AC-K33-ray-classes}
There are classes obtained from the actual preordinary
KLZ Rankin--Iwasawa classes,
\[
 z_{\mathfrak n}\in
 H^1\bigl(K(\mathfrak n),T\bigr)\otimes_E D_0^\vee ,
 \tag{RCN.8}
\]
with one fixed denominator in the chosen integral coefficient
lattices, such that, for every auxiliary prime
\(\mathfrak l\nmid\mathfrak n\),
\[
 \operatorname{Cor}_{K(\mathfrak n\mathfrak l)/K(\mathfrak n)}
        z_{\mathfrak n\mathfrak l}
   =P_{\mathfrak l}(\sigma_{\mathfrak l}^{-1})z_{\mathfrak n}.
 \tag{RCN.9}
\]
Here \(\sigma_{\mathfrak l}\) is arithmetic Frobenius in
\(\operatorname{Gal}(K(\mathfrak n)/K)\).
The formula includes split primes when their conjugate
already divides \(\mathfrak n\), and inert primes.
It uses the degree-two polynomial (RCN.3), not the
degree-four rational Rankin polynomial.
The two components indexed by \(D_0^\vee\) are compatible
classes; their nonzero singular images and independence
are not asserted.
\end{theorem}
\begin{proof}
\emph{1. The actual images and their integral normalization.}
Write \(R_N={}_c\mathrm{RI}^{[j]}_{1,N,0}\).
At each displayed tame level this notation denotes the compatible
family \(R_{Np^r}\) over the fixed base depth \(p^a\).
Its compatibility is the first case of KLZ Theorem 5.3.1,
since \(p\mid N\).  The away-prime identities below commute
with these transitions already on finite coefficient sheaves.
On the first factor use the moment
\(\mathrm{mom}^{k-j}\cdot\mathrm{id}\), then the fixed
eigenform quotient.  At higher tame level use ordinary
\(\operatorname{pr}_{1,*}\) to the fixed first level.
On the second factor use \(Q_{\mathfrak n}\) of
AC-K33-ray-norm.  The finite-coefficient Kunneth and
Hochschild--Serre maps of AC-K33-image are functorial
for all these correspondences.  They give actual classes
\[
 c_{\mathfrak n}\in
 H^1\!\left(\mathbb Q,
    \rho_f\otimes
     \rho_{\mathfrak n}^{\vee_{A_{\mathfrak n}}}(-j)\right)
           \otimes_{\mathcal O}D_0^\vee .
 \tag{RCN.10}
\]
The first moment is integral.  The fixed eigenform
projection loses only its fixed denominator, and
AC-K33-ray-norm supplies one fixed second-factor
denominator.  All source pushforwards, the edge maps
and the smoothing integer \(c\) are fixed or integral.
Thus the denominator is independent of \(\mathfrak n\).
No ordinary projector or inverse power of \(U_p'\)
is applied to the second factor.

Tensor-induction and duality identify the coefficient in
(RCN.10) with
\[
 \operatorname{Ind}_{G_K}^{G_{\mathbb Q}}
        \left(T\otimes_{\mathcal O}
                    A_{\mathfrak n}(\vartheta_{\mathfrak n}^{-1})\right).
 \tag{RCN.11}
\]
The inner module is the integral induced module from
\(G_{K(\mathfrak n)}\): in the group-element basis,
\(g\) sends \(t\otimes[\delta]\) to
\(gt\otimes[\bar g^{-1}\delta]\).
The two integral Shapiro maps therefore give (RCN.8).
For a ray quotient, the map \(A_{\mathfrak n'}\to
A_{\mathfrak n}\) sums the coefficients on its fibers.
In this induced-module description it is exactly the
trace map whose map on cohomology is corestriction.
This can also be checked from the finite-coset definition
of corestriction, term by term; no inverse group order
occurs.
For clarity, realize the induced module as functions
\(F:G_K\to T\) satisfying \(F(hx)=hF(x)\) for
\(h\in G_{K(\mathfrak n)}\), with right translation
\((gF)(x)=F(xg)\).  Writing \(F(x)=x t_{\bar x}\)
gives exactly the preceding coefficient action.  Multiplication
by \([\delta]\) sends \(F(x)\) to
\(x_\delta F(x_\delta^{-1}x)\); evaluating cocycles at
\(1\) identifies its cohomological action with
\(c(h)\mapsto x_\delta c(x_\delta^{-1}hx_\delta)\),
up to the explicit cocycle coboundary.  Thus multiplication
by the geometric Artin element \([\mathfrak l]\) acts
under Shapiro by \(\sigma_{\mathfrak l}^{-1}\),
with the usual Galois action on the extension field.

\emph{2. The three universal formulas.}
For a prime \(q\nmid N\), put \(s=q^{j+1}\).
At cyclotomic modulus one, the three source identities
and their symmetric version say
\[
\begin{split}
 (\operatorname{pr}_1,\operatorname{pr}_1)_*R_{Nq}
     &=(1-K_fK_g)R_N,\\
 (\operatorname{pr}_1,\operatorname{pr}_2)_*R_{Nq}
     &=(T_f-K_fT_g)R_N,\\
 (\operatorname{pr}_2,\operatorname{pr}_1)_*R_{Nq}
     &=(T_g-T_fK_g)R_N,\\
 (\operatorname{pr}_2,\operatorname{pr}_2)_*R_{Nq}
     &=s(1-K_fK_g)R_N .
\end{split}
 \tag{RCN.12}
\]
The first and second entries always refer to the fixed
first-factor and varying second-factor towers.
For \(q\mid N\), they instead read
\[
\begin{split}
 (\operatorname{pr}_1,\operatorname{pr}_1)_*R_{Nq}&=R_N,\\
 (\operatorname{pr}_1,\operatorname{pr}_2)_*R_{Nq}&=(U_f',1)R_N,\\
 (\operatorname{pr}_2,\operatorname{pr}_1)_*R_{Nq}&=(1,U_g')R_N,\\
 (\operatorname{pr}_2,\operatorname{pr}_2)_*R_{Nq}&=sR_N .
\end{split}
 \tag{RCN.13}
\]
These are equations before any eigenform or Artin quotient.
In particular they apply to the actual negative-moment
second map.  We now apply the already constructed
coefficient maps, using (RCN.4)--(RCN.5).

\emph{3. A split prime new at the rational level.}
Suppose \(q=\mathfrak l\bar{\mathfrak l}\), and neither
prime divides \(\mathfrak n\).
Over \(A_{\mathfrak n}\) let
\[
 a=u_{\mathfrak l},\qquad b=u_{\bar{\mathfrak l}},\qquad
 t=T_f,\quad d=d_f,\quad Q=q^{j+1}.
\]
Here \(T_g=a+b\), \(K_f=d/Q\), \(K_g=ab/Q\).
The actual coefficient norm is
\(\operatorname{pr}_{1,*}-aQ^{-1}\operatorname{pr}_{2,*}\).
The first factor is pushed down by \(\operatorname{pr}_{1,*}\).
The first two equations of (RCN.12) yield
\[
\begin{split}
 (1-K_fK_g)-\frac aQ(t-K_fT_g)
 &=1-\frac{at}{Q}
       +K_f\left(\frac{a(a+b)}Q-\frac{ab}Q\right)\\
 &=1-\frac{at}{Q}+\frac{a^2d}{Q^2}.
\end{split}
 \tag{RCN.14}
\]
This is the degree-two Euler factor, with no cancellation
of a class or an Euler factor.

\emph{4. A split prime whose conjugate is already present.}
Now let \(\bar{\mathfrak l}\mid\mathfrak n\) and
\(\mathfrak l\nmid\mathfrak n\).
The second local profile is the unramified tail \(v_a\),
so \(U_g'=a\).
Let \(F\) denote the first-factor quotient at level not
divisible by \(q\), and \(G\) this actual second-factor
quotient at level divisible by \(q\).
We first prove the relation
\[
 (F\operatorname{pr}_{2,*},G)_*R_{Nq}
       =a(F\operatorname{pr}_{1,*},G)_*R_{Nq},
 \qquad q\nmid N.
 \tag{RCN.15}
\]
Use \(R_{Nq}=(\operatorname{pr}_1,\operatorname{pr}_1)_*
R_{Nq^2}\), commute the two first-factor cyclic
degeneracies, and apply the third equation of
(RCN.13).  The result is the right side of (RCN.15),
because \(G U_g'=aG\).
This argument is over the full ramified group ring;
it does not first specialize the ramified character
to the trivial character.

Set \(F_1=F\operatorname{pr}_{1,*}\).
By (5.6.1a),
\[
 F_1U_f'=tF_1-K_fF\operatorname{pr}_{2,*}.
 \tag{RCN.16}
\]
Apply (RCN.13) to the norm
\(\operatorname{pr}_{1,*}-aQ^{-1}\operatorname{pr}_{2,*}\)
on the second factor, and use (RCN.15)--(RCN.16).
The resulting multiple of the lower class is
\[
 1-\frac{at}{Q}+\frac{a^2K_f}{Q}
       =1-\frac{at}{Q}+\frac{a^2d}{Q^2}.
 \tag{RCN.17}
\]
This proves the same factor in the conductor-drop case.

\emph{5. An inert prime.}
If \(\mathfrak l=(q)\), let \(u=u_{\mathfrak l}\).
At the lower level \(T_g=0\) and \(K_g=-u/Q\).
The actual second-factor norm is
\(\operatorname{pr}_{1,*}^2-uQ^{-2}
\operatorname{pr}_{2,*}^2\), with the first factor
pushed down by \(\operatorname{pr}_{1,*}^2\).
The first term on \(R_{Nq^2}\) is, by (RCN.13) then
(RCN.12),
\[
                   (1-K_fK_g)R_N.
 \tag{RCN.18}
\]
For the second term, first use
\((\operatorname{pr}_1,\operatorname{pr}_2)_*R_{Nq^2}
=(U_f',1)R_{Nq}\).
Then apply (5.6.1a) on the first factor and (RCN.12).
Its scalar is
\[
 t(t-K_fT_g)-K_f Q(1-K_fK_g)
       =t^2-K_fQ-K_f^2u .
 \tag{RCN.19}
\]
Combining (RCN.18)--(RCN.19) gives
\[
 1-\frac{(t^2-2d)u}{Q^2}
       +\frac{d^2u^2}{Q^4}.
 \tag{RCN.20}
\]
If \(\alpha,\beta\) are the two roots of the first
Hecke polynomial, then
\(t^2-2d=\alpha^2+\beta^2\).
Since \(\operatorname N\mathfrak l=q^2\), (RCN.20)
is precisely the degree-two Euler polynomial at
the inert \(K\)-prime.

\emph{6. Identification with (RCN.9).}
By (RCN.1)--(RCN.2), the geometric Frobenius roots of
\(T^\vee(1)\) at a split \(\mathfrak l\) are
\[
 \frac{\alpha_q(f)\,
       \chi^b(\operatorname{Frob}_{\mathfrak l}^{\rm geom})}
      {q^{j+1}},
 \qquad
 \frac{\beta_q(f)\,
       \chi^b(\operatorname{Frob}_{\mathfrak l}^{\rm geom})}
      {q^{j+1}}.
 \tag{RCN.21}
\]
At an inert prime, replace \(\alpha,\beta\) by their
squares and \(q^{j+1}\) by \(q^{2j+2}\).
Equations (RCN.14), (RCN.17), (RCN.20), followed by
the integral Shapiro comparison in Step 1, are therefore
exactly (RCN.9).
The coefficient maps commute on every finite ray
diagram by AC-K33-ray-norm, so these are compatible
classes for all squarefree auxiliary ideals.
Higher prime powers introduce no new \(p\)-primary
tame ray character and use the first identity in
(RCN.13).

The proof uses no assertion about the value of the
classes themselves.  Every displayed relation remains
valid if their projected images are zero.
\end{proof}

\begin{remark}[The remaining nonvanishing gate]
The actual class construction, the one fixed integral
denominator and the quadratic \(K\)-ray norm polynomial
are separate from an explicit reciprocity law.
To deduce the desired Bloch--Kato vanishing one still
needs the singular images of the two actual classes
to span the two-dimensional local target at \(269\),
with the precise source normalization.
The existence of the simultaneous quotient, the scalar
analytic interpolation of AC-K32, and the finite-level
geometric origin separately do not prove that determinant.
\end{remark}
